\chapter{Lire, chercher et démontrer}
\label{corriges:01}
\section*{Exercice 1.1 -- Vrai, faux, ou données insuffisantes}
\begin{enumerate}\item Vrai : $6k=3(2k)$.
\item Faux : $3$ est un multiple de $3$, pas de $6$.
\item Faux : deux triangles équilatéraux de côtés $2$ et $5$ ont tous leurs angles de $60\degre$.
\item Vrai : si $n$ est pair, le premier facteur est pair ; sinon $n+1$ est pair. Dans les deux cas, le produit est pair.\end{enumerate}

\section*{Exercice 1.2 -- La bonne implication}
La réciproque serait : « Si un quadrilatère a quatre côtés égaux, alors c'est un carré. » Elle est fausse : un losange avec des angles de $60\degre$ et $120\degre$ satisfait l'hypothèse mais pas la conclusion. Le mot attendu est \emph{losange}. Un carré est un losange particulier, qui possède aussi quatre angles droits.

\section*{Exercice 1.3 -- Démontrer avec une lettre}
Écrivons les entiers impairs $2p+1$ et $2q+1$. Leur somme vaut $2p+2q+2=2(p+q+1)$, multiple de $2$ puisque $p+q+1$ est entier. L'essai $3+7$ ne traite qu'une paire particulière ; le raisonnement avec $p,q$ traite toutes les paires.

\section*{Exercice 1.4 -- Réparer une rédaction}
Il manque une justification de l'angle droit. On peut disposer d'un codage ou d'une donnée explicite de perpendicularité ; le théorème direct devient alors applicable. Si l'on connaît exactement les trois côtés, on peut d'abord vérifier l'égalité des carrés et appliquer la réciproque. L'apparence du dessin n'est pas une hypothèse.
